% !PRINTED 216
\weberpage{216}
%<<CONT
$f(x) = 0$ setzt, simultane Invarianten von zwei Formen $n^{\text{ten}}$ und
$(n-2)^{\text{ten}}$ Grades, $f(z)$, $T(z)$ sind.

Eine solche ganze rationale und homogene symmetrische
Function $\nu^{\text{ter}}$ Ordnung

\[
P(y_1, y_2 \cdots y_n)
\]

ist wegen (1) offenbar eine homogene Function $\nu^{\text{ter}}$ Ordnung der
Variablen $t$. Sie ist ebenso eine homogene Function $\nu^{\text{ter}}$ Ordnung
 in den $a$, da der Ausdruck (1), wie er die $a$ explicite enth\"alt,
linear ist, und da die symmetrischen Functionen der $x$ nur
von den Verh\"altnissen $a_1 : a_0, a_2 : a_0, \ldots a_n : a_0$ abh\"angen. Bei dem
Ausdruck von $P$ durch die $a$ k\"onnte aber m\"oglicherweise eine
Potenz von $a_0$ im Nenner bleiben, und wir haben noch nachzuweisen,
dass dies nicht eintritt.

Setzen wir

\[
a_0^{\lambda} P(y_1, \ldots y_n) = K(t, a)
\]

und bestimmen die Potenz $a_0^{\lambda}$ so, dass $K(t, a)$ eine ganze Function
 der $a$ ist, aber f\"ur $a_0 = 0$ nicht mehr verschwindet, so
wird folgen, dass $\lambda = 0$ sein muss, wenn wir nachweisen, dass f\"ur
$a_0 = 0$ die s\"ammtlichen $y$ endliche Werthe behalten.

Wenn wir $a_0 = 0$ werden lassen, w\"ahrend die \"ubrigen $a$
unge\"andert bleiben, so wird eine der Wurzeln $x$, wie wir im
\S.\,40 gesehen haben, unendlich wachsen. Dass aber das zugeh\"orige
$y$ gleichwohl endlich bleibt, ergiebt sich aus der in Bezug auf
$t$ identischen Gleichung \S.\,68, (11)

\weq{2}{\frac{f(t)}{t-x} - \frac{1}{n}f'(t) = t^{n-2}F_0(x) + t^{n-3}F_1(x) + \cdots + F_{n-2}(x),}

aus der zu ersehen ist, dass, wenn $a_0 = 0$ und $x$ unendlich wird,
die Functionen

\[
F_0(x), \quad F_1(x) \cdots F_{n-2}(x)
\]

gleich den Co\"efficienten von $-\dfrac{1}{n}f'(t)$, also gleich

\[
0, \quad -\frac{n-1}{n}a_1 \ldots, \quad -\frac{1}{n}a_{n-1}
\]

werden. Es bleiben also auch die $y$ endlich, und $P(y_1, y_2 \cdots y_n)$
ist gleich einer \gesperrt{ganzen rationalen Invariante} vom Grade $\nu$
sowohl in den $t$ als in den $a$.

Wenn wir also die Gleichung f\"ur $y$ in der Form annehmen

\weq{3}{y^n + P_2 y^{n-2} + P_3 y^{n-3} + \cdots + P_n = 0,}

so ist $P_{\nu}$ eine solche Invariante $\nu^{\text{ten}}$ Grades.
